\documentclass[11pt,a4paper]{article}
\usepackage[margin=22mm]{geometry}
\usepackage{amsmath,amssymb,booktabs}
\input{glyphtounicode}
\pdfgentounicode=1
\usepackage[hidelinks]{hyperref}
\setlength{\parindent}{0pt}
\setlength{\parskip}{6pt}
\setlength{\abovedisplayskip}{5pt}
\setlength{\belowdisplayskip}{5pt}
\newcommand{\der}[1]{\frac{\mathrm d #1}{\mathrm d\alpha}}
\begin{document}
\begin{center}
{\Large Equation 17: stopping condition, time and discharged mass}
\end{center}

\textbf{Why the minima disappear.}
Equation 8 asks which condition first ends the present service interval.
Its minima select among head clearance, exhaustion of a sending allowance,
and the end of the timestep, across the participating lanes. Equation 17's
first line covers only the case where a head clearance wins that comparison.
It is not the derivative for every possible winner.

\begin{center}
\renewcommand{\arraystretch}{1.35}
\begin{tabular}{@{}ll@{}}
\toprule
Condition selected by equation 8 & Locally selected duration\\
\midrule
Head clearance in controlling lane $\ell$ & $\delta=q_\ell/r_\ell$\\
Sending allowance spent in controlling lane $\ell$
 & $\delta=S_\ell^{\rm rem}/r_\ell$\\
Timestep ends & $\delta=\Delta t-\sigma$\\
\bottomrule
\end{tabular}
\end{center}
If the same lane and condition continue to win under a small perturbation,
differentiate that selected expression. Other candidates do not determine
the interval locally. At a tie, this fixed-winner argument does not give
one universal ordinary derivative.

\textbf{Why the clearance-time derivative has two coefficients.}
Write $q,r$ for the controlling head's mass and rate. Clearing it obeys
the simple mass account
\[
 r\delta=q
 \quad\Longrightarrow\quad
 r\der{\delta}+\delta\der{r}=\der{q}.
\]
Solve for the change in required time:
\[
 \boxed{\der{\delta}
 =\underbrace{\frac1r\der{q}}_{\text{time needed for a mass change}}
  -\underbrace{\frac{\delta}{r}\der{r}}_{\text{time saved by a rate increase}},
 \qquad \frac{\delta}{r}=\frac{q}{r^2}.}
\]
The first coefficient converts an extra amount of traffic into extra
service time at the current rate. For the second term, a rate increase
would serve extra mass over the existing duration; dividing that extra
mass by the current rate gives how much earlier clearance occurs. These
are the two responses of the same mass account, not extra model weights.

\textbf{The discharge derivative has the complementary interpretation.}
For a given served stream, total discharge is the sum of rate times
duration over its service intervals:
\[
 y=\sum_a r_a\delta_a
 \quad\Longrightarrow\quad
 \der{y}=\sum_a\left(
 \underbrace{\delta_a\der{r_a}}_{\text{rate changes over existing time}}
 +\underbrace{r_a\der{\delta_a}}_{\text{time changes at existing rate}}
 \right).
\]
Thus the two parts of equation 17 connect through duration sensitivity.
The clearance formula supplies that sensitivity when a head determines
the endpoint; the other rows of the table govern the other stopping
conditions. For the final interval, remaining time changes when earlier
intervals change even though the timestep length itself is fixed.
\end{document}
